\documentclass{article}
\usepackage{mathtools} 
\usepackage{fontspec}
\usepackage[UTF8]{ctex}
\usepackage{amsthm}
\usepackage{mdframed}
\usepackage{xcolor}
\usepackage{amssymb}
\usepackage{amsmath}


% 定义新的带灰色背景的说明环境 zremark
\newmdtheoremenv[
  backgroundcolor=gray!10,
  % 边框与背景一致，边框线会消失
  linecolor=gray!10
]{zremark}{说明}

% 通用矩阵命令: \flexmatrix{矩阵名}{元素符号}{行数}{列数}
\newcommand{\flexmatrix}[4]{
  \[
  #1 = \begin{pmatrix}
    #2_{11}     & #2_{12}     & \cdots & #2_{1#4}   \\
    #2_{21}     & #2_{22}     & \cdots & #2_{2#4}   \\
    \vdots      & \vdots      & \ddots & \vdots     \\
    #2_{#31}    & #2_{#32}    & \cdots & #2_{#3#4}
  \end{pmatrix}
  \]
}

% 简化版命令（默认矩阵名为A，元素符号为a）: \quickmatrix{行数}{列数}
\newcommand{\quickmatrix}[2]{\flexmatrix{A}{a}{#1}{#2}}

\begin{document}
\title{总结 12}
\author{张志聪}
\maketitle

\begin{zremark}
  证明级数收敛的方法
\end{zremark}

\begin{itemize}
  \item 按照数列极限的方式处理。
  \item 把$S_n$看做整体，计算出$S_n$表达式。$\S 1$习题1-(5)
\end{itemize}

\begin{zremark}
  证明发散的方法
\end{zremark}

\begin{itemize}
  \item 级数收敛柯西准则的推论；
  \item $\lim \limits_{n \to \infty} u_n \neq 0$；
  \item 级数加括号后发散（定理12.4的推论）。
\end{itemize}

\begin{zremark}
  非负级数的性质
\end{zremark}

\begin{itemize}
  \item 部分和$\{S_n\}$单调递增。
\end{itemize}


\end{document}